% HD-F1-I02. Root draft; independent whole-proof review pending.
\section[HD-F1-I02: nine letters and one reflection]{HD-F1-I02: every rational nine-letter component with four nonreflection values and one reflection has period rank two}
\label{sec:f1-nine-letter-one-reflection-rank}

Retain the actual faithful monomial tower, character, corner
and polarization of HD-I01--HD-I21. Every component with
\[
 n=5,\qquad m=9,\qquad g(Y)=0,\qquad
 b=5,\qquad \nu=4,\qquad A=b-\nu=1
\]
has original period differential of rank two at every smooth
Hurwitz point. Here $d=3,4,6$ and $e\mid d$ are as in the
retained tower. A full $G(3,1,9)$ generating tuple below
establishes nonemptiness of this stratum.

\paragraph{1. The actual evaluation sequence.}
Let $Q$ contain the reduced uncolored cycles over the four
nonreflection values. The exact degree identity of I20 gives
\[
 \deg Q=(\nu-2)m+A-2n=9.
\]
For every scheme-theoretic fibre $D=h^*(z)$, of length nine,
the restriction sequence
\[
 0\longrightarrow\mathcal O_Y(Q-D)
 \longrightarrow\mathcal O_Y(Q)
 \longrightarrow\mathcal O_D(Q)\longrightarrow0
\]
has its first line bundle of degree zero on the actual
genus-zero curve. Thus its $H^0$ has dimension one and its
$H^1$ vanishes. The ten-dimensional space
$H^0(Y,\mathcal O_Y(Q))$ therefore evaluates surjectively
on every such fibre, including every nonreduced fibre.
Since the finite pushforward is locally free of rank nine,
these evaluations give the actual exact sequence
\[
 0\longrightarrow\mathcal K
 \longrightarrow H^0(Y,\mathcal O_Y(Q))\otimes\mathcal O
 \longrightarrow\mathcal V_Q=h_*\mathcal O_Y(Q)
 \longrightarrow0,
 \qquad \operatorname{rank}\mathcal K=1.
\]
Finite pushforward and Euler characteristic give
$\deg\mathcal V_Q=\deg Q+1-m=1$, so
$\deg\mathcal K=-1$. The degree classification of line
bundles on $\mathbf P^1$ (Stacks, tag~0BXJ,
Lemma~31.29.5) identifies this actual kernel with
$\mathcal O(-1)$.

Choose homogeneous coordinates $X,Y$ on the target and a
basis of the ten-dimensional section space. The kernel
inclusion is a column of linear forms, hence is
$X a+Y b$ for fixed vectors $a,b\in\mathbf C^{10}$.
The quotient is a vector bundle, so this column never
vanishes on the target. Consequently $a,b$ are linearly
independent: a dependent pair has a common zero on
$\mathbf P^1$. A constant change of basis therefore makes
the column $(X,Y,0,\ldots,0)^t$.
The actual map
\[
 \mathcal O^{\oplus2}\longrightarrow\mathcal O(1),
                 (u,v)\longmapsto -Yu+Xv
\]
has kernel the displayed column $\mathcal O(-1)$ and is
surjective at every point. The induced quotient isomorphism
gives
\[
             \mathcal V_Q\simeq
                   \mathcal O(1)\oplus\mathcal O^{\oplus8}.
\]
This is an isomorphism constructed from the actual evaluation
sequence and a constant change of basis.

\paragraph{2. The prescribed unit and the original period map.}
The literal constant function $1$ defines the given
inclusion $\mathcal O\to\mathcal V_Q$. It is nonzero in
every fibre by I21's local calculation: at a point in $Q$
with $w=z^\ell$ its coefficient in the generator $z^{-1}$
is $z$, whose class modulo $z^\ell$ vanishes exactly for
$\ell=1$. Vanishing on an entire fibre would thus make
every cycle uncolored and unramified, making that original
inertia the identity. Outside those four fibres the unit
is plainly nonzero. All original branch inertias are true.

In the actual isomorphism above this unit is $(s,c)$,
with $s\in H^0(\mathcal O(1))$ and $c\in\mathbf C^8$.
The vector $c$ is nonzero, since every section of
$\mathcal O(1)$ has a zero, and the unit has none as a
section of the pushforward bundle. A constant automorphism
of $\mathcal O^8$ sends $c$ to its first basis vector.
The triangular bundle automorphism
$(u,w)\mapsto(u-s w_1,w)$ then sends the unit to
$(0,1,0,\ldots,0)$. Its quotient is consequently
\[
 \mathcal W_Q=\mathcal V_Q/\mathcal O
             \simeq\mathcal O(1)\oplus\mathcal O^{\oplus7}.
\]
These explicit automorphisms retain the prescribed unit.
Using I21's actual tangent-fibre map gives
\[
 \ker d\mathcal F
   \simeq H^0\bigl(\mathcal W_Q(2-\nu)\bigr)
   =H^0\bigl(\mathcal O(-1)\oplus
                         \mathcal O(-2)^{\oplus7}\bigr)=0.
\]
The source has dimension $b-3=2$, so $d\mathcal F$ has
rank two at every smooth point.

Both actual primitive eigendifferential bundles on $Y$
have degree four. Their multiplication is onto
$H^0(\mathcal O(8))$: every degree-eight monomial is a
product of two degree-four monomials, as exhibited in
HD-F1-I01. Compose this actual multiplication with the
canonical quadratic trace of I20 and the polarized
isogeny of I02. The latter is the restriction of $q^*$,
with $q_*q^*=[|\ker\chi|]$ and the stated restricted
theta multiplier. These are exactly the original
period-map comparisons, so its differential also has
rank two everywhere on the stated smooth locus.

\paragraph{3. A connected full-group example.}
Use the actual matrices $C,D\in G(3,1,9)$ of HD-F1-I01.
Let
\[
 R=\operatorname{diag}(\zeta_3,1,\ldots,1),
 \qquad D'=DR^{-1},
 \qquad (C,C^{-1},D',R,D^{-1}).
\]
The matrix $D$ fixes the first coordinate, so this
definition removes exactly its color-one fixed cycle
there. The five entries have product one. Since
$D'R=D$, they generate the group generated by $C,D$,
which was proved by the explicit permutation and diagonal
commutators to be the full $G(3,1,9)$.

The moving ranks are respectively $8,8,5,1,6$.
Thus precisely the entry $R$ is a reflection, and the
other four are nonreflections. Their permutation
ramification sum is still $5+5+3+0+3=16$, giving the
genus-zero degree-nine intermediate quotient.
Their colored-cycle counts are $3,3,2,1,3$.
The color removed from $D$ is exactly the new color of
$R$, so both sums of primitive fractional colors remain
six. Both holomorphic multiplicities are therefore
$-1+6=5$.
The reduced uncolored divisor over the four nonreflection
values has counts $1,1,4,3$, summing to nine; it includes
unramified uncolored points. Riemann existence applied
to this actual full-group product-one tuple gives a
connected cover with its prescribed quotient tower.
This proves nonemptiness with all the stated invariants.

\paragraph{4. Framework and comparison scope.}
The original cover, finite character, local system,
middle-extension and Jacobian corner are retained. The
new morphisms are the displayed evaluation sequence,
its linear column quotient and the triangular automorphism
carrying the prescribed unit to a constant basis vector.
The actual polarized isogeny and quadratic trace identify
their tangent calculation with the original period map.
The Hilbert and cycle realizations, normalized Poincare
comparison and coherent $K_0$/Adams operations are the
same actual-tower maps recorded in F01--F04.
This computes the indicated individual component stratum
of T4-INDIVIDUAL. Its image has dimension two; all other
strata retain their previously stated gates.

For comparison, mark the actual twelve colored source
points and five target branch values. The same marked
Hurwitz diagram and unique cyclic-root comparison with
Filom's Thurston pullback used in HD-F1-I01 apply to
this degree-nine cover. The first assertion of his
Theorem~1.3 gives nonzero derivative because $12>9+2$.
The evaluation and prescribed-unit maps above compute
its exact rank two. The published nonvanishing precedent
is retained, with no priority assertion.
